Repetition-suppressing decoders usually fix a risk budget at every step, whereas what matters is the probability that a loop occurs anywhere in the sequence. We ask why a fixed per-step budget can waste KL divergence from the model, and by how much. We show that the loop probability falls by the visit-weighted sum of per-step hazard reductions, that to second order the cheapest allocation holds the price of risk (the Lagrange multiplier) constant, and that a fixed budget implies widely varying prices, its excess cost being one plus their squared coefficient of variation. In exactly solved finite-state systems a fixed price stays within 12.5\% of the minimum KL, whereas a fixed budget costs 1.3 to about 700 times the minimum, with and without infeasible budgets; the ratio tends to a finite limit (707 in the base system) as the reduction goes to zero. On a character-level language model a fixed price reached zero observed loops at 39 bits per generation, while every fixed budget below 1 tested cost at least 458 bits. These are observed counts at one temperature, not guarantees; KL measures distortion, not text quality; and about 98\% of each decoder's path KL lies beyond the cost of moving the loop probability.
